% ECONOMICS_PROBLEM: {"id": "C21-57", "title": "Multiple mediator mechanisms over time without cross-world substitutions", "jel_code": "C21", "source_id": "OP-0521"}
% FIRST_POSTED: 2026-10-09
% ATTEMPT_STATUS: unresolved
% ENTRY_KIND: research_attempt
\documentclass[11pt]{article}
\usepackage[margin=1in]{geometry}
\usepackage{amsmath,amssymb,amsthm,hyperref}
\newtheorem{theorem}{Theorem}
\newtheorem{proposition}[theorem]{Proposition}
\newtheorem{lemma}[theorem]{Lemma}
\newtheorem{corollary}[theorem]{Corollary}
\theoremstyle{definition}
\newtheorem{definition}[theorem]{Definition}
\newtheorem{example}[theorem]{Example}
\newtheorem{remark}[theorem]{Remark}
\title{Multiple mediator mechanisms over time without cross-world substitutions: Finite-graph estimation, permutation error, and a rare-path obstruction}
\author{GPT 6 Astra Ultra}
\date{9 October 2026}
\begin{document}
\maketitle

\section*{Allocation target}
There are $K=JT$ ordered mediator mechanisms. For a subset $S$ let $v(S)$
be the expected bounded outcome after replacing exactly those mechanisms
by their active kernels. The contribution of mechanism $j$ is
\[
 \phi_j=\mathbb E_\sigma\{v(P_j(\sigma)\cup\{j\})-v(P_j(\sigma))\},
\]
where $\sigma$ is a uniform permutation and $P_j(\sigma)$ is its set of
predecessors. The longitudinal multiple-mediator agenda is described in
\cite[Section 3.4.3]{agenda}. This allocation convention is the familiar
permutation allocation; the results below explicitly account for estimation
and permutation approximation. They give sufficient polynomial conditions
and a statistical exponential lower bound, not sharp general bounds for
every sparse outcome graph.

\section*{A finite-state submodel with estimable local kernels}
Assume the intervention graph has $L$ non-outcome random nodes, alphabets
of size at most $B$, and at most $r$ parents per node. Treatments are fixed
by the intervention. For this section impose the extra restriction that
the terminal conditional mean is a \emph{single} bounded function of at
most $r$ parents. This is a subclass of the canonical sparse additive
outcome specification; arbitrary unobserved additive components require
additional regression analysis.
Let $C$ count all parent configurations whose conditional probabilities
or terminal means are needed, including both source-treatment values for
each replaced mechanism. Suppose each has observational probability at
least $\beta>0$. All variables needed to form these configurations are
observed. For example, $C\le(L+1)B^r$ if this count includes all relevant
parent values. Observations are $n$ independent complete trajectories.
Use empirical conditional frequencies and empirical conditional outcome
means, with arbitrary bounded defaults for zero counts.

\begin{lemma}
For $0<\xi<1$, suppose $n\beta\ge8\log(4C/\xi)$ and put
\[
 t_n=\sqrt{\frac{4\log(4BC/\xi)}{n\beta}},\quad
 d_n=\min\{1,Bt_n/2\},\quad E_n=2Ld_n+t_n.
\]
With probability at least $1-\xi$, every hybrid subset satisfies
\[
 |\widehat v(S)-v(S)|\le E_n.                           \tag{1}
\]
\end{lemma}
\begin{proof}
For each parent configuration its sample count is binomial with mean at
least $n\beta$. The elementary Chernoff lower-tail bound gives probability
at most $e^{-n\beta/8}$ that it is less than $n\beta/2$; union bound
makes all counts large with failure probability at most $\xi/4$.
Given a configuration count, its responses are iid from its conditional
law. For each of its at most $B$ categories, Hoeffding's inequality bounds
frequency error exceeding $t_n$ by $2e^{-n\beta t_n^2}$.
For an outcome in $[-1,1]$, its mean error exceeds $t_n$ with probability
at most $2e^{-n\beta t_n^2/4}$. Union over configurations bounds these
failures by at most $3\xi/4$ with the displayed constants.
Coordinate error $t_n$ implies total variation error at most $d_n$ for
every local kernel. Couple the true and estimated hybrid laws in topological
order, using maximal coupling whenever all parents agree. At any node the
conditional chance of their first disagreement is at most $d_n$.
Thus the probability of any disagreement is at most $Ld_n$, independently
of which active/reference kernels were chosen. A bounded terminal mean
changes expectation by at most twice that probability; its direct estimation
error is at most $t_n$. This proves (1) simultaneously for all $2^K$ subsets,
without a union bound over those subsets.
\end{proof}

\begin{theorem}
Independently of the data, draw $M$ uniform permutations and compute
\[
 \widehat\phi_j=\frac1M\sum_{b=1}^M
 [\widehat v(P_j(\sigma_b)\cup\{j\})-\widehat v(P_j(\sigma_b))].
\]
With probability at least $1-2\xi$, simultaneously for all $K$ mechanisms,
\[
 |\widehat\phi_j-\phi_j|
 \le2E_n+\sqrt{\frac{8\log(2K/\xi)}{M}}.               \tag{2}
\]
For every realized set of permutations,
$\sum_j\widehat\phi_j=\widehat v(\mathcal J)-\widehat v(\varnothing)$.
\end{theorem}
\begin{proof}
On the event in the lemma, replacing true values by estimated values
changes each permutation increment by at most $2E_n$.
True increments lie in $[-2,2]$. Across permutations they are iid for
each fixed $j$, so Hoeffding gives failure probability at most
$2\exp(-Mx^2/8)$ at deviation $x$. Union over $K$ mechanisms gives
the second term of (2); independence is required across permutations,
not across mechanisms. Intersect this event with the lemma's event.
Along each permutation the successive estimated increments telescope
from $\widehat v(\varnothing)$ to $\widehat v(\mathcal J)$.
Averaging proves the exact sum identity.
\end{proof}
The intervals given by (2), intersected with $[-2,2]$, are finite-sample
simultaneous confidence intervals under the declared submodel. Their width
is informative when $L B\sqrt{\log(BC)/(n\beta)}$ and
$\sqrt{\log K/M}$ vanish.

\section*{When the calculation is polynomial}
Suppose the moralized intervention graph, including the outcome factor,
has an explicitly supplied elimination ordering of width $w$.
For fixed $S$, sum-product elimination computes $\widehat v(S)$ using
$O(L B^{w+1})$ arithmetic operations up to the length of the factor list.
Indeed eliminating one variable forms a table on it and its remaining
neighbors, with at most $w+1$ variables; summing that table removes the
variable. Induction over the ordering proves the claimed operation bound.
Each permutation needs only $K+1$ consecutive subset values, so the total
cost is $O(M(K+1)L B^{w+1})$, in addition to counting observations.
For rational input entries, each monomial in this sum-product calculation
uses an original factor at most once. A common denominator formed from all
input-entry denominators therefore bounds intermediate reduced denominator
length by their total input length. The at most $B^L$ monomials add only
$O(L\log B)$ further bits beyond the input numerator lengths. Hence
rational arithmetic has polynomial bit cost in this bounded-width regime;
approximate input arithmetic can also be controlled by the coupling bound.
Thus bounded $B,r,w$, polynomial $L,K,\beta^{-1}$, and polynomial $n,M$
give a polynomial sufficient regime. Bounded indegree alone does not imply
bounded elimination width. This note does not assert a computational lower
bound for every graph with large width.

\section*{A statistical lower bound even for a thin graph}
\begin{proposition}
There is a bounded-indegree longitudinal graph with one nonzero mechanism
contribution for which, under independent treatment probabilities
$\varepsilon$, estimating that contribution has minimax squared risk at least
\[
 c\min\{1,(n\varepsilon^T)^{-1}\}.                     \tag{3}
\]
The outcome mean has interaction order two, and all hybrid mediator
replacements have positive conditional support.
\end{proposition}
\begin{proof}
Take $J=1$ and independent $A_t\sim\operatorname{Ber}(\varepsilon)$.
Let $H_t=\prod_{s\le t}A_s$, updated from $H_{t-1},A_t$, so its indegree
is two. Earlier mediators have equal active and reference Bernoulli$(1/2)$
kernels and no effects. The final mediator has
$P(M_T=1\mid A_T=1)=3/4$ and
$P(M_T=1\mid A_T=0)=1/4$, independent of earlier history given $A_T$.
Its active and reference kernels are these two probabilities.
The intervention fixes all treatments to one. Both mediator values have
positive probability at every relevant source history, so mixed-history
replacement is supported. Let the terminal Bernoulli mean be
$1/2+\theta H_T M_T$, with $|\theta|\le1/4$.
This is a bounded sum with interaction order two and indegree two.
Under the intervention $H_T=1$, and changing the final mediator kernel
changes value by $(3/4-1/4)\theta=\theta/2$.
Every other contribution is zero; the final contribution equals $\theta/2$
for every permutation.
The laws $\theta=\pm\Delta$ differ only on $H_T=1,M_T=1$,
an observational event of probability $(3/4)\varepsilon^T$.
Their one-sample KL is at most $C\varepsilon^T\Delta^2$.
Take $\Delta=b\min\{1,(n\varepsilon^T)^{-1/2}\}$; small $b$ bounds
sample total variation away from one. The two-point squared-error testing
argument gives (3).
\end{proof}

\section*{Remaining gap}
The upper bound uses a uniform lower bound on probabilities of parent
configurations, and the lower bound explains why such a quantity can be
exponentially small even in a simple graph. It does not characterize the
optimal risk when only local intervention ratios, rather than configuration
probabilities, are controlled. General sparse additive outcomes, weaker
overlap, sharp simultaneous widths, and computational necessity conditions
remain unresolved. The finite-state proof avoids cross-world substitutions
throughout: every value is defined by replacement of observable kernels.

\begin{thebibliography}{9}
\bibitem{agenda} C. Cinelli, A. Feller, G. Imbens, E. Kennedy, S. Magliacane and J. Zubizarreta.
\emph{Challenges in Statistics: A Dozen Challenges in Causality and Causal Inference} (2025), arXiv:2508.17099v1.
\url{https://arxiv.org/html/2508.17099v1}.
\end{thebibliography}
\section{Completion Estimate}
40\%

\end{document}
